\begin{statement}{C.2.UN.10}{Proved}{A contracted irrational horizontal section is elliptic without ambient Betti assumptions}
Let $f:Y\to B$ be a proper holomorphic map from a
smooth complex threefold to a smooth curve, with
connected fibres and smooth general fibre a complex
two-torus. Let $mS$ be an irreducible nonnormal full
fibre. Write $h=\eta\pi:Z\to T\to S$ for the
normalization and its minimal resolution. Suppose
that $Z$ is ruled over a curve of positive genus,
that its effective adjunction divisor is
$D=2C+V$, with $C$ a section and $V$ vertical,
and that $\pi$ contracts $C$. Then the ruling
base has genus one.
\end{statement}
\begin{proof}
We first supply the comparisons that previously
came from the global Betti assumptions. On a
small full fibre tube $U$, the rank-one torsion-free
sheaf $f_*K_U$ is a line bundle. A generator has
no horizontal zeros, since its restriction to a
smooth torus is a nonzero canonical form. Its
zero divisor is consequently $kS$ for some
$k\ge0$, after shrinking away other special
fibres. Thus the evaluation map gives the actual
isomorphism
\[
 K_U\simeq\OO_U(kS),\qquad
 \omega_S\simeq\OO_U((k+1)S)|_S.                \tag{C.UNCR.1}
\]
The divisor $mS$ is principal, so both these line
bundles are torsion. The map to the smooth curve
is flat, and the holomorphic Euler characteristic
of its full fibres is constant. The filtration
of $\OO_{mS}$ has quotients
$\OO_S(-jS)$ for $0\le j<m$; each is a torsion
line-bundle twist of $\OO_S$. Analytic
Riemann--Roch gives the same Euler characteristic
for these twists, since their rational Chern
characters are one. It follows that
\[
            0=\chi(\OO_{mS})=m\,\chi(\OO_S).
 \tag{C.UNCR.2}
\]
The exact normalization and resolution duality
maps therefore give $K_Z+D=h^*\omega_S$, with
the same effective divisor and torsion right-hand
side as in C.2.CR.1--C.2.CR.2.

Here are the parts of that argument needed here,
with no transgression-class assumption. Finite
duality embeds $\omega_T$ into the torsion line
$\eta^*\omega_S$, with its sections vanishing on
the nonempty conductor. A nonzero section of a
torsion line on a connected compact normal space
is nowhere zero: taking a trivializing power
makes it a nonzero constant. Consequently
$H^2(T,\OO_T)=0$. The coherent Leray sequence
then bounds $\operatorname{length}R^1\pi_*\OO_Z$
by $h^1(Z,\OO_Z)=g$, whereas formal functions
and restriction to the contracted section $C$
give length at least $h^1(C,\OO_C)=g$ at its
image point $p$. Thus equality holds at $p$ and
there are no further contributions.

At $q=\eta(p)$, fix a local generator of $\omega_S$.
Finite duality identifies the branch conductor
ideal $I_p$ with $\omega_{T,p}$, embedded as an
actual ideal summand in $\OO_{S,q}$. Proper
duality and the residue sequence of $C$ give
exactly
\[
 I_p=(\pi_*\omega_Z(C))_p,\qquad
 I_p/(\pi_*\omega_Z)_p\simeq H^0(C,\omega_C).
 \tag{C.UNCR.3}
\]
The equality of lengths follows from the preceding
formal-functions computation. In the fixed local
dualizing trivialization these sheaves are
$\OO_Z(-C-V)$ and $\OO_Z(-2C-V)$.
The maximal ideal of $\OO_{S,q}$ vanishes on $C$;
its square supplies an additional order along
$C$ without losing any prescribed order along
$V$. Hence
$I_p\cap\mathfrak m_q^2\subset(\pi_*\omega_Z)_p$.
The embedding dimension three and Nakayama
argument of C.2.CR.2 now give $g\le2$. When
$g=2$ they give the actual conductor ideal
$(x,y)$, a unique normalization point $p$ above
$q$, and generic transverse delta invariant
one. Thus the full conductor-image curve is
smooth and reduced at $q$, and the generic
transversal is a node or a cusp.

Every positive-dimensional critical component
of the reduced surface is part of that conductor
curve; after shrinking, its full reduced critical
locus is exactly the smooth curve just obtained.
The node case is a product node by C.2.UN.9.
Its normalization has two smooth points above
$q$, contradicting the unique point $p$.
The cusp case is a product cusp by C.2.CR.4.
Its explicit normalization
$(u,z)\mapsto(X,Y,z)=(u^3,-u^2,z)$, after
adjusting the harmless sign in $X^2+Y^3$,
is smooth. This contradicts a minimal
normalization resolution contracting a
positive-genus curve to $p$. Both genus-two
possibilities are excluded, and $g=1$ follows.
\end{proof}


% C.2.UN.42 from part_C2_nonlinear_hopf_occurrence.tex

